\documentclass[a4paper, 11pt]{report}

\usepackage[utf8]{inputenc}                                                      
\usepackage[french]{babel}                                                       
\usepackage[T1]{fontenc}  
\usepackage[pdftex]{graphicx}                                                    
\usepackage{url}                                                                 
\usepackage{longtable}
\usepackage[bookmarks, colorlinks=false, pdfborder={0 0 0},
	pdftitle={Medicagenda}, 
	pdfauthor={Knop Florian, Kriwin Paul, Placentino Simon},
	pdfsubject={Medicagenda},           
pdfkeywords={UML, ISO/IEC 19505-1:2012, Medicagenda, Projet, analyse, ESI}]{hyperref}   

\newcommand{\HRule}{\rule{\linewidth}{0.5mm}}    

\begin{document}
\begin{titlepage}
	\begin{center}

		\includegraphics[keepaspectratio=true,width=0.20\textwidth]{../ressources/logo}\\[1cm]

		\textsc{\LARGE H.E.B. Ecole Superieur d'Informatique}\\[1.5cm]

		\textsc{\Large Laboratoire d'analyse : Projet d'analyse}\\[0.5cm]
		\textsc{\Large Use Case Specifications - UC 1106 - Envoyer une
		notification}\\[0.5cm]

		\HRule \\[0.4cm]
		{\huge \bfseries Medicagenda \\[0.4cm]}
		\HRule \\[1.5cm]

		\noindent
		\begin{minipage}[t]{0.4\textwidth}
			\begin{flushleft} \large
				\emph{Auteurs:}\\
				Florian \textsc{Knop} \href{mailto:39310@heb.be}{39310@heb.be}\\
				Paul \textsc{Kriwin} \href{mailto:39171@heb.be}{39171@heb.be}\\
				Simon \textsc{Placentino} \href{mailto:39631@heb.be}{39631@heb.be}\
			\end{flushleft}
		\end{minipage}%
		\begin{minipage}[t]{0.4\textwidth}
			\begin{flushright} \large
				\emph{Titulaire du cours:} \\
				Mr.~Nicolas \textsc{Pettiaux}
				\href{mailto:npettiaux@heb.be}{npettiaux@heb.be}
			\end{flushright}
		\end{minipage}

		\vfill

		{\large \today}

	\end{center}
	\clearpage\null\newpage
\end{titlepage}

\tableofcontents

\chapter{UC Specification : Envoyer une notification}

\section{Introdcution}

\subsection{Objectifs de ce document}

Ce document présente les spécifications fonctionnelles du UC1106 - Envoyer une notification.

\subsection{Domaine de définition de ce document}

\section{Définition du Use Case}

\subsection{Identifiant et nom}

UC1106 - Envoyer une notification.

\subsection{Brève description}

Avant un rendez-vous, le système envoie un rappel au patient pour le prévenir d'un rendez-vous.
Le patient peut choisir le temps avant le rendez-vous pour envoyer le rappel ainsi que le moyen
de communication (e-mail, sms, etc.).

\section{Flux}

\subsection{Flux de base}

Cet UC est déclenché automatiquement un certain nombre d'heure avant un rendez-vous. Quand un patient ou un mèdecin
définissent un rendez-vous, un envoi de notification est programmé. L'heure est défini grâce aux paramètres utilisateur, de même
pour le type d'envoi. Le patient peut annuler un rendez-vous, si le rendez-vous est annulé alors que la notification est déjà
programmé, la notification est annulée et une nouvelle notification confirmant l'annulation est envoyé directement au 
patient.

\subsection{Flux alternatif}

\subsubsection{Interruption du UC}

Si le système subit une panne, il faudra contacter l'administrateur système.
Celui-ci lancera un UC pour relancer les notifications non envoyées avec un message d'excuse accompagnant ces
notifications.


\section{Acteurs, mode, etc.}

\subsection{Acteurs}

Le SIA est l'acteur principale, c'est à dire qu'en cas de panne, il faudra s'adresser au 
responsable du système.

\subsection{Mode}

Automatisé unitaire.

\subsection{Evènement déclencheur}

Après que le patient ou le mèdecin ait créé un rendez-vous avec l'UC1104. 
Un envoi de notification est programmé à une certaine heure avant l'heure du rendez-vous selon les 
paramètres utilisateurs du patient.

\section{Pré-conditions}

Le patient doit avoir un rendez-vous de prévu.

\section{Post-conditions}

Aucune.

\section{Point d'inclusions et d'extensions}

SO. (Sans Objets).
\section{Règles de gestion}

\subsection{Classes concernées}

La classe Rendez-vous et Compte Patient sont concernées en lecture seulement.

\subsection{Validation des encodages de données}

Aucune données n'est transmise de cette manière au système.

\subsection{Validation des données en sortie}

Aucune données transmises par le système ne nécessite une telle vérification.

\section{Diagramme d'activité}

Cet UC ne possède pas de diagramme d'activité, il est du ressort du système d'agir en fonctions
des choix de l'utilisateur.

\section{Interface utilisateur}

Etant un UC automatisé, cet UC ne possède pas d'interface graphique.


\end{document}
